\documentclass[../main.tex]{subfiles}
\begin{document}

\chapter{高度なキャッシュ}
\lessoninfo{
  video={Lesson 14，動画 1--40},
  textbook={H\&P \cite{hennessy2017} §2.3，§B.1，§B.3},
  keywords={パイプライン化キャッシュ，VIPT キャッシュ，エイリアシング，ウェイ予測，NMRU，疑似 LRU，競合性ミス，プリフェッチ，ループ交換，ミスアンダーミス対応，MSHR，キャッシュ階層，大域ミス率，包含性},
}

第12章では基本的なキャッシュを述べ，平均メモリアクセス時間でそれを測った．
第13章では，仮想記憶をもつシステムでどのメモリアクセスにも必要となる
アドレス変換を加えた．この章では，実際のキャッシュを高速にする技法を
まとめる．技法は，平均メモリアクセス時間のどの項を下げるかによって
整理する．すなわち，ヒット時間（パイプライン化キャッシュ，TLB に対する
キャッシュの置き方，ウェイ予測，安価な置換アルゴリズム），ミス率
（ブロックの大型化，プリフェッチ，ループ交換），そしてミスペナルティ
（重なったミスの処理と，何よりも多段のキャッシュ階層）である．

\section{キャッシュ性能の改善}
\label{sec:l14-improving}

改善すべき量は，第12章の平均メモリアクセス時間（AMAT）である．
\source{Lesson 14，動画 1--2；H\&P §2.3．}
\begin{equation}
  \text{AMAT} = \text{ヒット時間} + \text{ミス率} \times \text{ミスペナルティ} .
  \label{eq:l14-amat}
\end{equation}
3つの項のそれぞれに固有の攻め方があり，\cref{tab:l14-techniques} には
この章の技法を，対象とする項ごとに挙げる．代償を伴わないものはほとんど
ない．一方の項を下げる技法は，しばしば別の項を上げる---キャッシュを
小さくすればヒット時間は短くなるがミス率は高くなる---ので，AMAT 全体が
減る場合にのみ価値がある．

\begin{table}[htbp]
  \centering\small
  \caption{この章の技法を，\cref{eq:l14-amat} のどの項を下げるかで
    分類したもの．}
  \label{tab:l14-techniques}
  \begin{tabular}{@{}lll@{}}
    \toprule
    項 & 技法 & 節 \\
    \midrule
    ヒット時間   & パイプライン化キャッシュ，VIPT キャッシュ，    & \labelcref{sec:l14-pipelined}--\labelcref{sec:l14-replacement} \\
                 & ウェイ予測，NMRU 置換と PLRU 置換             & \\
    ミス率       & ブロックの大型化，プリフェッチ，ループ交換     & \labelcref{sec:l14-three-cs}--\labelcref{sec:l14-loop-interchange} \\
    ミスペナルティ & ミスアンダーミス対応，キャッシュ階層         & \labelcref{sec:l14-overlap}--\labelcref{sec:l14-inclusion} \\
    \bottomrule
  \end{tabular}
\end{table}

\section{パイプライン化キャッシュ}
\label{sec:l14-pipelined}

ヒット時間を下げる最も直接的な方法は，キャッシュを小さく単純に---たとえば
ダイレクトマップに---保つことである．デコードすべきラインが少なく，比較
すべきタグも少なくて済む．\source{Lesson 14，動画 3--4；H\&P §2.3．}
しかし小さなキャッシュはミスが多く，ダイレクトマップキャッシュは競合を
多く受けるので，ヒットで節約した時間はミスで容易に失われる．この節と
続く4つの節の技法は，ミス率を犠牲にすることなく，ヒット時間またはその
代償を減らす．

第12章で述べたように，L1 のヒットには1〜3クロックサイクルかかる．
\caution{ベンダが公表する L1 のレイテンシは，配列アクセス単体ではなく，
ロードから値が使えるまで（load-to-use）の値である．Skylake は単純な
アドレスで 4 サイクル，インデックス付きで 5 サイクルとする．Intel
最適化マニュアル．}数
サイクルを要する探索が，完了するまでキャッシュを占有するとすれば，どの
メモリアクセスも直前のアクセスを待たなければならない．アクセスから見た
ヒット時間は自身の探索時間にこの待ち時間を加えたものとなり，連続する
ロードとストアはプロセッサを止めてしまう．探索が1サイクルに収まる
キャッシュにはこの問題はない．2〜3サイクルかかる L1 キャッシュは，通常
パイプライン化される．

\begin{definition}[パイプライン化キャッシュ]
  \term[ぱいぷらいんかきゃっしゅ]{パイプライン化キャッシュ}[pipelined cache]%
  とは，探索をラッチで区切られた段に分割し，1回の探索には数クロック
  サイクルかかるが，毎サイクル新しい探索を開始できるようにしたキャッシュ
  である．
\end{definition}

セットアソシアティブキャッシュでは，探索は自然に3つの手順に分かれ，その
それぞれが段になりうる（\cref{fig:l14-pipelined}）．
\begin{enumerate}
  \item インデックスがセットを選び，そのセットのラインの有効ビットとタグを
    読み出す．
  \item タグをアドレスのタグと比較して，ヒットかミスかを決めるとともに
    一致したウェイを特定し，データの読み出しを開始する．
  \item データの読み出しが完了し，一致したウェイのデータが選択され，
    ブロックオフセットがブロックから要求されたバイトを取り出して引き渡す．
\end{enumerate}
ミスは第2段の終わりに判明する．\margin{第2段でタグと並んですべての
ウェイのデータを読むことが，探索を速くしている．大きな L2 や L3 はまず
タグを読み，一致したウェイのデータだけを後から読む．遅いが消費電力は
はるかに小さい．}パイプライン化は，1回の探索にかかる
サイクル数を減らすものではない．そのサイクル数がアクセスを発行できる
頻度を制限しないようにするものである．あるアクセスがタグを比較している
間に，次のアクセスはすでに自身のセットを読み出しており，これは第3章のパイプラインで
命令が重なるのとまったく同じである．そしてクロック周期は，探索全体では
なく最も遅い段だけを収めればよい．

\begin{figure}[htbp]
  \centering
  % Three consecutive accesses to a set-associative cache whose lookup is
% split into three pipeline stages: a new access starts in every cycle.
\begin{tikzpicture}[x=1.95cm,y=1.0cm, font=\footnotesize\sffamily,
  st/.style={draw, minimum width=1.85cm, minimum height=0.9cm, inner sep=1pt,
    align=center, font=\scriptsize\sffamily},
  s1/.style={st, fill=ln-box},
  s2/.style={st, fill=ln-accent!15},
  s3/.style={st, fill=ln-tip!20}]
  \foreach \c in {1,...,5} { \node at (\c,0.75) {\c}; }
  \node[anchor=east] at (0.45,0.75) {\iflangja{サイクル}{cycle}};
  \foreach \i/\n in {0/1,1/2,2/3} {
    \node[anchor=east] at (0.45,-\i) {\iflangja{アクセス}{access}~\n};
    \node[s1] at (\i+1,-\i) {\iflangja{セット選択\\V・タグ読出し}{select set,\\read V and tags}};
    \node[s2] at (\i+2,-\i) {\iflangja{タグ比較，\\データ読出し開始}{compare tags,\\start data read}};
    \node[s3] at (\i+3,-\i) {\iflangja{読出し完了，\\ウェイ選択，\\オフセットで抽出}{finish read,\\select way,\\offset $\to$ bytes}};
  }
\end{tikzpicture}

  \caption{探索を3段に分割したキャッシュへの，連続する3回のアクセス．
    各アクセスには3サイクルかかるが，新しいアクセスが毎サイクル開始される．}
  \label{fig:l14-pipelined}
\end{figure}

\begin{example}
  \label{ex:l14-pipelined}
  ある L1 キャッシュの探索論理の遅延は \SI{0.75}{\nano\second} であり，
  プロセッサの他の段は \SI{0.3}{\nano\second} のクロック周期を許す．
  パイプライン化しなければ，プロセッサ全体のクロック周期が
  \SI{0.75}{\nano\second} 以上に伸びるか，各アクセスが3サイクルにわたって
  キャッシュを占有してメモリアクセスが3サイクルに1回しか開始できなくなる
  かである．ラッチのオーバーヘッドの余地を残して約 \SI{0.25}{\nano\second}
  の3段に分割すれば，探索は \SI{0.3}{\nano\second} のクロックに収まる．
  データはアクセス開始の3サイクル後に到着し，アクセスは毎サイクル開始
  できる．\margin{パイプライン化が与えるのは毎サイクル1アクセスである．
  それ以上はポートとバンクによる．Skylake の L1 データキャッシュは1
  サイクルに2つのロードと1つのストアを処理する．Intel 最適化マニュアル．}
\end{example}

\section{アドレス変換とヒット時間}
\label{sec:l14-translation}

仮想記憶（第13章）のもとでは，プロセッサは仮想アドレスを発行するが，
主記憶のデータは物理アドレスで位置づけられる．\source{Lesson 14，動画
5--7；H\&P §B.3．}キャッシュの探索は，変換の後に行うことも，前に行うことも，
一部を変換と並列に行うこともでき（\cref{fig:l14-translation}），この選択が
ヒット時間に影響する．

\begin{figure}[htbp]
  \centering
  % Where address translation sits relative to the cache lookup:
% (a) physically accessed, (b) virtually accessed, (c) VIPT cache.
\begin{tikzpicture}[x=1cm,y=1cm, font=\footnotesize\sffamily, >={Stealth[length=1.8mm]},
  blk/.style={draw, fill=ln-box, minimum height=0.6cm, inner sep=2pt},
  cch/.style={draw, fill=ln-accent!15, minimum height=0.7cm, minimum width=2.2cm, inner sep=2pt},
  fld/.style={draw, fill=white, minimum height=0.5cm, inner sep=1pt, font=\scriptsize\sffamily},
  lab/.style={font=\scriptsize\sffamily, text=ln-muted, align=left},
  ttl/.style={font=\footnotesize\sffamily, align=center}]
  % ---------------- (a) physically accessed
  \node (va) at (1.6,0) {\iflangja{仮想アドレス}{virtual address}};
  \node[blk, minimum width=1.4cm] (tlb) at (1.6,-1.2) {TLB};
  \node[cch] (c) at (1.6,-2.7) {\iflangja{キャッシュ}{cache}};
  \node (d) at (1.6,-3.8) {\iflangja{データ}{data}};
  \draw[->] (va) -- (tlb);
  \draw[->] (tlb) -- node[lab, right] {\iflangja{物理\\アドレス}{physical\\address}} (c);
  \draw[->] (c) -- (d);
  \node[ttl] at (1.6,-4.55) {\iflangja{(a) 物理アドレスキャッシュ}{(a) physically accessed}};
  % ---------------- (b) virtually accessed
  \node (vb) at (5.2,0) {\iflangja{仮想アドレス}{virtual address}};
  \node[cch] (cb) at (4.9,-2.7) {\iflangja{キャッシュ}{cache}};
  \node (db) at (4.9,-3.8) {\iflangja{データ}{data}};
  \node[blk, dashed, fill=white, minimum width=1.2cm] (tb) at (6.3,-1.2) {TLB};
  \node[lab, anchor=north] at (6.3,-1.5) {\iflangja{変換は\\ミス時のみ}{translation\\only on a miss}};
  \draw[->] (4.9,-0.25) -- (cb);
  \draw[->, dashed] (6.3,-0.25) -- (tb);
  \draw[->] (cb) -- (db);
  \node[ttl] at (5.3,-4.55) {\iflangja{(b) 仮想アドレスキャッシュ}{(b) virtually accessed}};
  % ---------------- (c) VIPT
  \node[fld, minimum width=1.3cm] (vpn) at (9.05,0) {VPN};
  \node[fld, minimum width=1.3cm, anchor=west] (idx) at (vpn.east) {\iflangja{インデックス}{index}};
  \node[lab, anchor=south] at (9.7,0.25) {\iflangja{仮想アドレス}{virtual address}};
  \node[blk, minimum width=1.2cm] (tc) at (9.05,-1.2) {TLB};
  \node[cch, minimum width=2.6cm] (cc) at (9.7,-2.7) {\iflangja{キャッシュ}{cache}};
  \node (dc) at (9.7,-3.8) {\iflangja{データ}{data}};
  \draw[->] (vpn) -- (tc);
  \draw[->] (tc) -- node[lab, left] {\iflangja{フレーム\\番号}{frame\\number}} (tc.south |- cc.north);
  \draw[->] (idx.south) -- (idx.south |- cc.north);
  \draw[->] (cc) -- (dc);
  \node[ttl] at (9.7,-4.55) {(c) VIPT};
\end{tikzpicture}

  \caption{TLB とキャッシュ探索を組み合わせる3つの方法．(b) では変換は
    ミス時にだけ必要なので，TLB はヒットの経路に入らない（権限の確認の
    ためには依然として参照される）．(c) では，インデックスがすでにセットを
    選んでいる間に TLB が仮想ページ番号を変換する．}
  \label{fig:l14-translation}
\end{figure}

\subsection{物理アドレスキャッシュ}

物理アドレスで探索されるキャッシュを\term[ぶつりあどれすきゃっしゅ]{物理アドレスキャッシュ}[physically accessed cache]%
という．どのアクセスもまず TLB で仮想アドレスを変換し，得られた物理
アドレスだけが探索に使われる（\cref{fig:l14-translation}a）．したがって
ヒット時間には変換が含まれる．
\begin{equation}
  \text{ヒット時間} = \text{TLB の探索時間} + \text{キャッシュの探索時間} .
  \label{eq:l14-pa-hit}
\end{equation}
TLB は小さく高速だが，その遅延はヒットを含むすべてのアクセスで支払われ，
まさに我々が減らそうとしている量に上乗せされる．\margin{TLB は追いつく
ために小さくなければならない．Skylake の第1レベルデータ TLB は 4~KB
ページ用に 64 エントリであり，その後ろに第2レベルの 1536 エントリが
ある．Intel 最適化マニュアル．}

\subsection{仮想アドレスキャッシュ}

\begin{definition}[仮想アドレスキャッシュ]
  \term[かそうあどれすきゃっしゅ]{仮想アドレスキャッシュ}[virtually accessed cache]%
  は，インデックスもタグも仮想アドレスから取る．アドレスが変換されるのは，
  ブロックを主記憶から取得しなければならないミスのときだけである．
\end{definition}

ヒット時には TLB は経路に入らず，ヒット時間はキャッシュの探索時間そのもの
である（\cref{fig:l14-translation}b）．この利点には問題が伴う．
\begin{itemize}
  \item \textbf{権限．} ページの TLB エントリは，そのページを読めるか，
    書けるか，実行できるかも記録しており，権限はすべてのアクセスで確認
    しなければならない．そのため実際には，ヒット時にも TLB が参照される．
  \item \textbf{コンテキストスイッチ．} 各プロセスは自分の仮想アドレス
    空間をもつので，同じ仮想アドレスが異なるプロセスでは異なるデータを
    指す．オペレーティングシステムがプロセッサをあるプロセスから別の
    プロセスへ切り替える（コンテキストスイッチ）と，キャッシュ内の
    ブロックは古いプロセスの仮想アドレスで識別されており，新しい
    プロセスがそれらを見つけてはならない．すべての有効ビットを落として
    キャッシュをフラッシュするか，すべてのタグにアドレス空間の識別子を
    付け加えるかである（H\&P §B.3）．コンテキストスイッチは毎秒何度も
    起こり，フラッシュする方式では，次に走るプロセスが毎回，空の
    キャッシュとミスの多発から始めることになる．
  \item \textbf{エイリアシング．} これは \cref{sec:l14-aliasing} の主題で
    ある．
\end{itemize}

\subsection{仮想インデックス物理タグキャッシュ}

両者を組み合わせると，短いヒット時間を保ちながらフラッシュを避けられる．

\begin{definition}[VIPT キャッシュ]
  \term[かそういんでっくすぶつりたぐきゃっしゅ]{仮想インデックス物理タグキャッシュ}[virtually indexed physically tagged cache]%
  \index{VIPT|see{仮想インデックス物理タグキャッシュ}}（VIPT キャッシュ）は，
  インデックスを仮想アドレスから，タグを物理アドレスから取る．
\end{definition}

仮想アドレスのインデックスビットが直ちにセットを選ぶ
（\cref{fig:l14-vipt}）．そのセットの有効ビットとタグが読み出されている
間に，TLB が仮想ページ番号を物理フレーム番号に変換し，物理的な値である
セットのタグが TLB からのフレーム番号と比較される．
\begin{itemize}
  \item ヒット時間は本質的にキャッシュの探索時間である．TLB の探索は
    セットの選択と並列に走り，タグが比較される時点までにフレーム番号を
    届ければよく，小さな TLB なら通常これができる．
  \item コンテキストスイッチでキャッシュをフラッシュする必要がない．タグは
    ブロックの由来する物理メモリを表しており，どのプロセスが走っていても
    それは正しいままである．
  \item いずれにせよ TLB はすべてのアクセスで参照されるので，権限の確認に
    追加の費用はかからない．
\end{itemize}

\begin{figure}[htbp]
  \centering
  % Lookup in a VIPT cache: 32-bit addresses, 4 KB pages, 64-byte blocks,
% 64 sets of 2 lines (8 KB).  The index comes from the page offset of the
% virtual address; the TLB translates the VPN in parallel, and the frame
% number is compared with the (physical) tags of the set; a line hits only if
% its valid bit is set.
\begin{tikzpicture}[x=1cm,y=1cm, font=\footnotesize\sffamily, >={Stealth[length=1.8mm]}]
  % --- virtual address fields
  %% The Japanese index label is wider than its box in the English layout,
  %% so in Japanese the VPN/index boundary (and the page-offset bracket that
  %% starts there) sits further left; the TLB arrow at x=1.3 and the index
  %% line at x=6.2 stay inside their boxes.
  \iflangja{%
    \draw[fill=ln-accent!15] (0.6,0) rectangle (4.0,0.55);
    \node at (2.3,0.275) {VPN\ (31\,\dots\,12)};
    \draw[fill=ln-box] (4.0,0) rectangle (7.0,0.55);
    \node at (5.5,0.275) {インデックス\ (11\,\dots\,6)};
  }{%
    \draw[fill=ln-accent!15] (0.6,0) rectangle (4.6,0.55);
    \node at (2.6,0.275) {VPN\ (31\,\dots\,12)};
    \draw[fill=ln-box] (4.6,0) rectangle (7.0,0.55);
    \node at (5.8,0.275) {index\ (11\,\dots\,6)};
  }
  \draw[fill=white] (7.0,0) rectangle (9.4,0.55);
  \node at (8.2,0.275) {\iflangja{オフセット}{offset}\ (5\,\dots\,0)};
  \node[anchor=south west, font=\sffamily\scriptsize, text=ln-muted, inner sep=1pt] at (0.6,0.6) {\iflangja{仮想アドレス}{virtual address}};
  % page offset bracket
  \iflangja{%
    \draw[ln-muted] (4.02,0.62) -- (4.02,0.78) -- (9.38,0.78) -- (9.38,0.62);
    \node[anchor=south, font=\sffamily\scriptsize, text=ln-muted] at (6.7,0.78)
      {ページオフセット (11\,\dots\,0)};
  }{%
    \draw[ln-muted] (4.62,0.62) -- (4.62,0.78) -- (9.38,0.78) -- (9.38,0.62);
    \node[anchor=south, font=\sffamily\scriptsize, text=ln-muted] at (7.0,0.78)
      {page offset (11\,\dots\,0)};
  }
  % --- TLB
  \node[draw, fill=ln-box, minimum width=1.3cm, minimum height=0.6cm] (tlb) at (1.3,-1.3) {TLB};
  \draw[->] (1.3,0) -- (tlb.north);
  % --- two ways, four drawn sets
  \foreach \x in {2.4,6.6} {
    \fill[ln-accent!15] (\x,-1.55) rectangle (\x+3.4,-2.0);
    \node[font=\sffamily\scriptsize, text=ln-muted, anchor=south] at (\x+0.2,-1.1) {V};
    \node[font=\sffamily\scriptsize, text=ln-muted, anchor=south] at (\x+1.1,-1.1) {\iflangja{タグ}{tag}};
    \node[font=\sffamily\scriptsize, text=ln-muted, anchor=south] at (\x+2.6,-1.1) {\iflangja{データ}{data}};
    \foreach \i in {0,...,3} {
      \draw (\x,-1.1-0.45*\i) rectangle (\x+0.4,-1.55-0.45*\i);
      \draw (\x+0.4,-1.1-0.45*\i) rectangle (\x+1.8,-1.55-0.45*\i);
      \draw (\x+1.8,-1.1-0.45*\i) rectangle (\x+3.4,-1.55-0.45*\i);
    }
  }
  \node[font=\sffamily\scriptsize, text=ln-muted, anchor=north] at (5.0,-2.95) {\iflangja{ウェイ 0}{way 0}};
  \node[font=\sffamily\scriptsize, text=ln-muted, anchor=north] at (9.2,-2.95) {\iflangja{ウェイ 1}{way 1}};
  \foreach \i/\s in {0/0,1/1,2/{\dots},3/63} {
    \node[font=\sffamily\scriptsize, text=ln-muted, anchor=west] at (10.05,-1.325-0.45*\i) {\s};
  }
  \node[font=\sffamily\scriptsize, text=ln-muted, anchor=south west] at (10.0,-1.1) {\iflangja{セット}{set}};
  % --- index selects a set in both ways
  \draw (6.2,0) -- (6.2,-1.775);
  \fill (6.2,-1.775) circle (1.2pt);
  \draw[->] (6.2,-1.775) -- (5.8,-1.775);
  \draw[->] (6.2,-1.775) -- (6.6,-1.775);
  % --- comparators against the frame number from the TLB
  \node[draw, circle, fill=white, inner sep=1pt, minimum size=0.5cm] (c0) at (3.5,-3.7) {$=$};
  \node[draw, circle, fill=white, inner sep=1pt, minimum size=0.5cm] (c1) at (7.7,-3.7) {$=$};
  \draw[->] (3.5,-2.9) -- (c0.north);
  \draw[->] (7.7,-2.9) -- (c1.north);
  % valid bits enable the comparators
  \draw[->] (2.6,-2.9) -- (c0.north west);
  \draw[->] (6.8,-2.9) -- (c1.north west);
  \draw (tlb.south) -- (1.3,-4.9) -- (7.7,-4.9);
  \draw[->] (7.7,-4.9) -- (c1.south);
  \fill (1.3,-3.7) circle (1.2pt);
  \draw[->] (1.3,-3.7) -- (c0.west);
  \node[anchor=east, align=right, font=\sffamily\scriptsize, text=ln-muted] at (1.2,-2.7)
    {\iflangja{物理\\フレーム\\番号}{physical\\frame\\number}};
  \node[draw, fill=ln-box, minimum width=0.7cm, minimum height=0.5cm, inner sep=1pt, font=\sffamily\scriptsize] (or) at (5.6,-3.7) {OR};
  \draw[->] (c0.east) -- (or.west);
  \draw[->] (c1.west) -- (or.east);
  \draw[->] (or.south) -- (5.6,-4.3) node[below] {\iflangja{ヒット}{hit}};
\end{tikzpicture}

  \caption{4~KB のページ，64 バイトのブロック，64 セット（8~KB）をもつ 2
    ウェイ VIPT キャッシュにおける探索．インデックスは仮想アドレスの
    ページオフセットから取られる．タグは物理フレーム番号であり，TLB の
    出力と比較される．ラインは，有効ビットが立っていてタグが一致すれば
    ヒットする．}
  \label{fig:l14-vipt}
\end{figure}

\section{エイリアシング}
\label{sec:l14-aliasing}

仮想記憶では，同じ物理メモリが複数の仮想アドレスに現れうる．
\source{Lesson 14，動画 8--11；H\&P §B.3．}

\begin{definition}[エイリアシング]
  同じアドレス空間内の，あるいは異なるアドレス空間の2つ以上の仮想アドレスが
  同じ物理アドレスに対応づけられる状況を
  \term[えいりあしんぐ]{エイリアシング}[aliasing]という．このとき，それらの
  仮想アドレスは互いにエイリアスである．
\end{definition}

エイリアシングはよく起こる．メモリを共有するプロセスは，通常は異なる仮想
アドレスで同じフレームを見るし，1つのプロセスが1つのフレームを2つの
アドレスに対応づけることもある．仮想アドレスキャッシュでは，2つの
エイリアスは単に2つの異なるアドレスである．一般にインデックスもタグも
異なるので，キャッシュは同じ物理ブロックの写しをエイリアスごとに1つずつ
保持でき，写し同士を結びつけるものは何もない．

\begin{example}
  \label{ex:l14-alias}
  仮想アドレス V1 と V2 は物理アドレス P のエイリアスであり，仮想アドレス
  キャッシュでは両者のインデックスが異なる．P のメモリには値~5 が入って
  いる．\cref{tab:l14-alias} は4回のアクセスを追う．最後のアクセスは V2 を
  通して~5 を読むが，同じメモリには V1 を通して~7 が書き込まれている．
  ライトスルー方式であれば主記憶は~7 を保持するが，それでも V2 用の古い
  写しが返される．
\end{example}

\begin{table}[htbp]
  \centering\small
  \caption{ライトバック方式の仮想アドレスキャッシュにおける2つのエイリアス
    （\cref{ex:l14-alias}）．各アクセスの後の，2つの写しと主記憶の内容．}
  \label{tab:l14-alias}
  \begin{tabular}{@{}clllll@{}}
    \toprule
    手順 & アクセス      & 結果            & V1 の写し & V2 の写し & P のメモリ \\
    \midrule
    1    & V1 を読む     & ミス            & 5         & ---       & 5 \\
    2    & V2 を読む     & ミス            & 5         & 5         & 5 \\
    3    & V1 に 7 を書く & ヒット          & 7         & 5         & 5 \\
    4    & V2 を読む     & ヒット，5 を返す & 7         & 5         & 5 \\
    \bottomrule
  \end{tabular}
\end{table}

正しく動作するためには，どの書き込みも同じ物理ブロックのキャッシュ内の
他のすべての写しを見つけ出し，更新するか無効化しなければならないが，
キャッシュにはそれらを効率的に探す手段がない．したがってエイリアシングは
仮想アドレスキャッシュのもう1つの問題である．\sidenote{オペレーティング
システムの側で2つ目の写しを防ぐこともできる．ページカラーリングでは，
エイリアスをインデックスビットの一致する仮想アドレスにしか割り当てないので，
エイリアスは同じセットを選ぶ．代償は，どのフレームをページに与えるかの
自由度である（H\&P §B.3）．}

\subsection{VIPT キャッシュにおけるエイリアシング}

VIPT キャッシュではタグが物理的なので，ブロックの2つの写しは同じタグを
もつことになる．それでも両者は異なるセットに入りうる．インデックスが仮想
アドレスから取られるからである．それが起こりうるかどうかは，インデックス
ビットがどこにあるかで決まる．変換は仮想ページ番号をフレーム番号に
置き換え，ページオフセットはそのまま残す．したがってエイリアスは，仮想
ページ番号こそ異なるが，ページオフセットは常に等しい．
\begin{itemize}
  \item インデックスが仮想ページ番号のビットを含む場合，2つのエイリアスは
    異なるセットを選びうるので，キャッシュは2つの写しを保持でき，
    \cref{ex:l14-alias} の問題が残る．
  \item インデックスがページオフセットのビットだけからなる場合，
    エイリアスは常に同じセットを選ぶ．そこには自分のブロックの物理タグを
    もつラインが1つあるだけなので，写しはたかだか1つしか存在せず，
    エイリアシングは無害である．このとき仮想アドレスのインデックスビットは
    物理アドレスのインデックスビットでもあり，キャッシュは物理アドレス
    キャッシュとまったく同じに，ただしより高速に振る舞う．
\end{itemize}

\subsection{エイリアシングの回避}

したがって VIPT キャッシュは，ブロックオフセットとインデックスがページ
オフセットに収まるとき，エイリアシングから自由である．
\begin{equation}
  \text{オフセットビット数} + \text{インデックスビット数} \;\le\; \text{ページオフセットビット数} .
  \label{eq:l14-vipt-bits}
\end{equation}
左辺は $\log_2(\text{セット数} \times \text{ブロックサイズ})$，右辺は
$\log_2(\text{ページサイズ})$ なので，この条件は
$\text{セット数} \times \text{ブロックサイズ} \le \text{ページサイズ}$ と
読める．両辺に連想度~$N$ を掛けると，キャッシュサイズの上界が得られる．
\begin{equation}
  \text{キャッシュサイズ} \;\le\; N \times \text{ページサイズ} .
  \label{eq:l14-vipt-size}
\end{equation}
ブロックサイズは消えた．VIPT キャッシュを制限するのは連想度とページ
サイズだけである．\tip{バイトではなくビットで数える．4~KB ページで変換を
受けないのはビット 0--11 だけなので，オフセットビット数 $+$ インデックス
ビット数 $\le 12$ である．}

\begin{example}
  \label{ex:l14-vipt-size}
  4~KB のページではページオフセットは 12 ビットである．64 バイトの
  ブロック（オフセット 6 ビット）をもつ VIPT キャッシュは，インデックスを
  最大 6 ビット，すなわち 64 セットまでしか使えない．2 ウェイならこれは
  \cref{fig:l14-vipt} の 8~KB のキャッシュであり，4 ウェイなら 16~KB で
  ある．16~KB の 2 ウェイキャッシュは 128 セットを必要とし，その7番目の
  インデックスビットは仮想アドレスのビット~12，すなわち仮想ページ番号の
  最下位ビットとなるので，エイリアスが異なるセットに置かれうる．
\end{example}

\subsection{実際の VIPT キャッシュ}

ページサイズはアーキテクチャとオペレーティングシステムが定めるもので，
典型的には 4~KB である．\cref{eq:l14-vipt-size} により，VIPT であり，かつ
エイリアシングから自由であろうとする L1 キャッシュは，連想度を上げること
によってしか大きくできない．ダイレクトマップなら 4~KB が限界であり，
32~KB のキャッシュには 8 ウェイが必要である．Intel の Core~2，Nehalem，
Sandy Bridge，Haswell はいずれも 32~KB の 8 ウェイ L1 データキャッシュを
もち，ちょうどこの上界にある．\margin{ページを大きくすれば上界も上がる．
Apple の M シリーズは 16~KB ページを用い，M1 の高性能コアは 128~KB の L1
データキャッシュをもつ．Apple silicon CPU 最適化ガイド．}したがって実際の L1 キャッシュは，ヒット
時間だけでは正当化できないほど連想度が高く，そのことが次の2つの節の技法を
重要にしている．

\begin{keypoint}
  VIPT キャッシュは，TLB が並列に変換する間に仮想アドレスでインデックス
  し，物理タグを比較する．ヒット時間は仮想アドレスキャッシュとほぼ同じで，
  コンテキストスイッチでのフラッシュを必要とせず，
  $\text{キャッシュサイズ} \le N \times \text{ページサイズ}$ である限り
  エイリアシングの影響を受けない．
\end{keypoint}

\section{ウェイ予測}
\label{sec:l14-way-prediction}

連想度はヒット時間とミス率を逆向きに動かす．\source{Lesson 14，動画
12--15；H\&P §2.3．}セット当たりのウェイが多いほど競合は減ってミス率は
下がり---VIPT キャッシュでは許されるサイズも大きくなる
（\cref{eq:l14-vipt-size}）---が，探索ではより多くのタグを比較し，より
多くのラインから選択しなければならないので，ヒット時間は長くなる．
ダイレクトマップキャッシュはヒット時間が最も短く，ミス率が最も高い．
目標は，セットアソシアティブキャッシュのミス率を，ダイレクトマップ
キャッシュのヒット時間で得ることである．

\begin{definition}[ウェイ予測]
  \term[うぇいよそく]{ウェイ予測}[way prediction]を行うセットアソシアティブ
  キャッシュは，選ばれたセットのどのラインが要求されたブロックを保持して
  いるかを推測し，まずそのラインだけを調べる．そのタグが一致すれば，
  アクセスはダイレクトマップキャッシュのヒット時間で完了する．そうでなければ，
  通常のセットアソシアティブキャッシュと同様にセットの他のラインを調べる．
\end{definition}

サイズ~$C$ の $N$ ウェイキャッシュの最初の確認はセット当たり1つのラインを
調べるので，同じセットをもつサイズ $C/N$ のダイレクトマップキャッシュの
ように振る舞い，そのようなキャッシュのヒット率が，推測の当たる頻度を
見積もる．推測がセット内で最も最近アクセスされたラインであれば，この
見積もりは厳密である．そのラインはセット内で最も最近アクセスされた
ブロックを保持しており，それはまさにサイズ $C/N$ のダイレクトマップ
キャッシュが保持するブロックだからである．推測が外れれば，素早い確認に
加えて通常の探索の分も費やし，ミスもまた通常の探索の後にしか判明しない．
予測されたラインがヒットするアクセスの割合を~$p$，1つのラインを調べる
時間を~$t_1$，セット全体の探索時間を~$t_N$ とすると，平均探索時間は次で
与えられる．
\begin{equation}
  t = p\, t_1 + (1-p)(t_1 + t_N) = t_1 + (1-p)\, t_N .
  \label{eq:l14-way-prediction}
\end{equation}

\begin{example}
  \label{ex:l14-way-prediction}
  ある L1 キャッシュの後ろには L2 キャッシュがあり，L1 のミスペナルティは
  \SI{20}{\nano\second} である．3つの L1 の設計を比較する．4~KB の
  ダイレクトマップキャッシュは，ウェイ予測が最初に調べる 8 ウェイ
  キャッシュの部分集合なので，予測されたラインは全アクセスの
  $p = \SI{80}{\percent}$ でヒットし，残りの \SI{20}{\percent}（他の
  ラインでのヒットとすべてのミス）は $0.8 + 1.5 = \SI{2.3}{\nano\second}$
  を要する．
  \begin{center}\small
    \begin{tabular}{@{}lccl@{}}
      \toprule
      L1 キャッシュ & 探索時間 & ミス率 & AMAT \\
      \midrule
      4~KB ダイレクトマップ     & \SI{0.8}{\nano\second} & \SI{20}{\percent} & $0.8 + 0.20\times 20 = \SI{4.8}{\nano\second}$ \\
      32~KB 8 ウェイ            & \SI{1.5}{\nano\second} & \SI{5}{\percent}  & $1.5 + 0.05\times 20 = \SI{2.5}{\nano\second}$ \\
      32~KB 8 ウェイ，予測あり  & 0.8 または \SI{2.3}{\nano\second} & \SI{5}{\percent} & $1.1 + 0.05\times 20 = \SI{2.1}{\nano\second}$ \\
      \bottomrule
    \end{tabular}
  \end{center}
  \cref{eq:l14-way-prediction} により，予測ありの平均探索時間は
  $0.8 + 0.2\times 1.5 = \SI{1.1}{\nano\second}$ である．ウェイ予測は
  8 ウェイキャッシュのミス率を保ったまま，その平均探索時間を
  ダイレクトマップキャッシュのそれに近づける．
\end{example}

\section{置換アルゴリズムとヒット時間}
\label{sec:l14-replacement}

置換アルゴリズムはミス率だけでなくヒット時間にも影響する．
\source{Lesson 14，動画 16--19；H\&P §B.1．}第12章のアルゴリズムのうち
2つが両極端をなす．
\begin{description}
  \item[ランダム置換] ヒット時には何もせず，ミス時に犠牲を無作為に選ぶ．
    ヒット時間には何も加えないが，これから使われるブロックを追い出して
    しまうことがあるので，ミス率は高くなる．
  \item[LRU] はミス率が低いが，各ラインに $\log_2 N$ ビットのカウンタを
    保持し，ヒットであれミスであれ1回のアクセスがセットの $N$ 本すべての
    ラインのカウンタを変えうる．すべてのアクセスでのこの活動は論理と
    エネルギーとヒット時間を費やし，連想度が高いほどその度合いは大きい．
\end{description}
以下のアルゴリズムは，はるかに少ない管理情報で LRU のミス率に近づく．

\subsection{NMRU}

\begin{definition}[NMRU]
  \term{NMRU}[not most recently used]方式は，各セットについてどのラインが
  最も最近アクセスされたかだけを記録し，ミス時にはそれ以外のラインの1つを
  無作為に選んで追い出す．
\end{definition}

時間的局所性により，最も最近使われたブロックは次にアクセスされる可能性が
最も高い．それを追い出すことは，ランダム置換がなしうる最も有害な選択で
あり，NMRU はまさにその選択を除外する．費用はセット当たり $\log_2 N$
ビットのポインタ1つであり，1回のアクセスが変えるのはたかだかこのポインタ
だけである．8 ウェイキャッシュには，LRU の $8 \times 3 = 24$ ビットでは
なく，セット当たり 3 ビットで足りる．NMRU は LRU よりいくらか多くミスする．
その犠牲が，最も最近使われたものよりわずかに古いだけのブロックでありうる
からである．

\subsection{疑似 LRU}

\begin{definition}[PLRU]
  \term[ぎじ LRU]{疑似 LRU}[pseudo-LRU]\index{PLRU|see{疑似 LRU}}（PLRU）
  方式は，ライン当たり1ビットを初期値~0 で保持する．アクセスは自分の
  ラインのビットを立てる．これによってセット内でまだ~0 であった最後の
  ビットが立つと，他のすべてのラインのビットが落とされる．ミスは，ビットが~0 の
  ラインを無作為に選んで追い出す．
\end{definition}

ラインのビットは，そのラインが最後のリセット以降にアクセスされたかどうかを
示す．\caution{他の文献で「疑似 LRU」というときは，木構造の変種を指すことが
多い．セット当たり $N-1$ ビットが二分木をたどって犠牲を選ぶ．ここでの
ライン当たり1ビットは別の近似である．}最も最近使われたラインのビットは常に立っているので，PLRU の候補は
常に NMRU の候補に含まれ，PLRU は NMRU と同等以上である．LRU にどれだけ
近いかは，いくつのビットが立っているかによる．リセットの直後には，NMRU と
同じく，最も最近使われたラインだけが除外される．$N-1$ 個のビットが立って
いるときは，唯一の候補は最後のアクセスが最も昔であるラインであり，これは
まさに LRU の選択である．費用はセット当たり $N$ ビットで，1回のアクセスは
リセットの場合を除いて1ビットを立てるだけである．

\begin{example}
  \label{ex:l14-replacement}
  ある 4 ウェイのセットが，ライン 0--3 にブロック A，B，C，D を保持して
  いる．最も最近使われたのは D であり，PLRU のビットは D へのアクセスで
  最後にリセットされた．\cref{tab:l14-replacement} は，3つのアルゴリズムの
  もとでのこのセットを追う．E でのミスでは，LRU は A を追い出さなければ
  ならず，NMRU は A，C，D のいずれかを，PLRU は A か C を追い出しうる．
  ここでは3つとも A を追い出す．次のアクセス F もミスすれば，LRU は C を
  追い出し，NMRU は B，C，D のいずれかを追い出しうる．PLRU はビットが~0 の
  ラインがライン~2 だけなので，やはり C を追い出す．
\end{example}

\begin{table}[htbp]
  \centering\small
  \caption{1つの 4 ウェイセットについて LRU，NMRU，PLRU が保持する状態
    （\cref{ex:l14-replacement}）．LRU の順は最も長く使われていない
    ブロックから最も最近使われたブロックまで，NMRU は最も最近使われた
    ラインの記録，PLRU はライン 0--3 のビットである．}
  \label{tab:l14-replacement}
  \begin{tabular}{@{}llcccc@{}}
    \toprule
    アクセス & 結果 & ライン 0--3 & LRU の順 & NMRU & PLRU のビット \\
    \midrule
    （初期状態） &      & A B C D & A B C D & 3 & 0 0 0 1 \\
    B         & ヒット & A B C D & A C D B & 1 & 0 1 0 1 \\
    D         & ヒット & A B C D & A C B D & 3 & 0 1 0 1 \\
    B         & ヒット & A B C D & A C D B & 1 & 0 1 0 1 \\
    E         & ミス   & E B C D & C D B E & 0 & 1 1 0 1 \\
    \bottomrule
  \end{tabular}
\end{table}

\begin{keypoint}
  すべてのアクセスで行われる作業は，すべてのヒットを長くする．ウェイ予測は
  他のラインより先に1つのラインだけを調べ，NMRU と PLRU は LRU のカウンタを，
  セット当たり1つのポインタ（NMRU）またはライン当たり1ビット（PLRU）に
  置き換える．いずれも，連想度や LRU によるミス率の利点の大半を，ヒット
  時間の代償のわずかな割合で保つ．
\end{keypoint}

\section{ミスの原因}
\label{sec:l14-three-cs}

ミス率を下げるには，まずミスがなぜ起こるかを知ることから始める．
\source{Lesson 14，動画 20；H\&P §B.3；\cite{hill1989}．}ミスは3種類に
分類され，しばしば 3C と呼ばれる．

\begin{definition}[初期参照ミス，容量性ミス，競合性ミス]
  \term[しょきさんしょうみす]{初期参照ミス}[compulsory miss]とは，ある
  ブロックへの最初のアクセスで生じるミスである．これはどんなキャッシュでも，
  無限に大きなキャッシュでさえ生じる．

  \term[ようりょうせいみす]{容量性ミス}[capacity miss]とは，同じサイズの
  フルアソシアティブキャッシュでも生じるであろうミスである．キャッシュが
  プログラムの使っているすべてのブロックを保持できず，場所が足りないために
  そのブロックが追い出されていたのである．

  \term[きょうごうせいみす]{競合性ミス}[conflict miss]とは，連想度が限られて
  いるためだけに生じるミスである．同じサイズのフルアソシアティブキャッシュ
  ならまだそのブロックを保持しているが，実際のキャッシュでは，同じセットを
  争う別のブロックによって追い出されていた．
\end{definition}

3つの個数は，同じアクセス列を3つのキャッシュでシミュレートすれば測定
できる．無限のキャッシュのミスが初期参照ミスであり，実際のサイズの
フルアソシアティブキャッシュがそれに容量性ミスを加え，実際のキャッシュが
さらに競合性ミスを加える．

\begin{example}
  \label{ex:l14-three-cs}
  あるプログラムについて，無限のキャッシュは \num{1500} 回，32~KB の
  フルアソシアティブキャッシュは \num{6000} 回，実際の 4 ウェイ 32~KB
  キャッシュは \num{7800} 回ミスする．\num{7800} 回のミスのうち
  \num{1500} 回が初期参照ミス，$\num{6000} - \num{1500} = \num{4500}$ 回が
  容量性ミス，$\num{7800} - \num{6000} = \num{1800}$ 回が競合性ミスで
  ある．
\end{example}

この分類は対処法を示す．キャッシュを大きくすれば容量性ミスが減り，連想度を
上げるか置換アルゴリズムを改善すれば競合性ミスが減るが，いずれもヒット
時間を長くする．\margin{マルチプロセッサでは4つ目の C が加わる．他のコアの
書き込みがブロックを無効化して生じるコヒーレンスミスである（第19章）．}以下の節では，ヒットを長くすることなくミス率を下げる技法を
述べる．

\section{キャッシュブロックの大型化}
\label{sec:l14-block-size}

与えられたサイズのキャッシュでは，ブロックサイズを大きくすることが，
ミス率に相反する2通りの影響を与える．\source{Lesson 14，動画 21--22；
H\&P §B.3．}
\begin{itemize}
  \item 空間的局所性が良い場合，アクセスされたアドレスの周辺のデータは
    すぐに使われる．ブロックが大きければ1回のミスでそれらを多くキャッシュに
    持ち込め，そうでなければミスしていた後続のアクセスがヒットするように
    なる．
  \item 空間的局所性が乏しい場合，大きなブロックの大半は一度も使われない．
    キャッシュが保持できるブロックの数が減り，したがって実際に必要とされる
    ブロックも減って，容量性ミスと競合性ミスが増える．
\end{itemize}
そのためブロックサイズを大きくしていくと，ミス率はまず下がり，その後再び
上がる（\cref{fig:l14-block-size}）．最小値は，キャッシュが大きいほど
大きなブロックサイズに現れる．たとえば 4~KB のキャッシュでは約 64 バイト，
256~KB のキャッシュでは約 256 バイトである．大きなキャッシュはブロックが
大きくても多くのブロックを収める余地があるので，使われないデータに費やす
場所の害が小さい．

\begin{figure}[htbp]
  \centering
  % Miss rate as a function of the block size for a 4 KB and a 256 KB cache
% (shape only, no data).  The best block sizes follow the lecture: about
% 64 bytes for the small and about 256 bytes for the large cache.
\begin{tikzpicture}[x=1cm,y=1cm, font=\footnotesize\sffamily, >={Stealth[length=1.8mm]}]
  \draw[->] (0,0) -- (7.6,0);
  \node[font=\footnotesize\sffamily] at (3.8,-0.75) {\iflangja{ブロックサイズ（バイト）}{block size (bytes)}};
  \draw[->] (0,0) -- (0,4.4) node[above, inner sep=2pt] {\iflangja{ミス率}{miss rate}};
  \foreach \x/\l in {0.8/16,2.0/32,3.2/64,4.4/128,5.6/256,6.8/512} {
    \draw (\x,0) -- (\x,-0.08) node[below, font=\scriptsize\sffamily] {\l};
  }
  \draw[thick, ln-accent] plot[smooth, tension=0.7] coordinates
    {(0.5,3.95) (0.8,3.6) (2.0,3.0) (3.2,2.82) (4.4,3.1) (5.6,3.65) (6.8,4.2)};
  \draw[thick, ln-tip] plot[smooth, tension=0.7] coordinates
    {(0.5,2.25) (0.8,1.95) (2.0,1.35) (3.2,1.0) (4.4,0.82) (5.6,0.74) (6.8,0.9) (7.2,1.02)};
  \fill[ln-accent] (3.2,2.82) circle (1.6pt);
  \fill[ln-tip] (5.6,0.74) circle (1.6pt);
  \node[anchor=west, text=ln-accent] at (7.0,4.2) {\iflangja{4~KB キャッシュ}{4~KB cache}};
  \node[anchor=west, text=ln-tip] at (7.3,1.02) {\iflangja{256~KB キャッシュ}{256~KB cache}};
  \node[font=\scriptsize\sffamily, text=ln-muted, align=center] at (1.9,2.3)
    {\iflangja{近傍へのアクセスが\\ヒットする}{accesses to neighbours\\now hit}};
  \node[font=\scriptsize\sffamily, text=ln-muted, align=center] at (5.3,2.3)
    {\iflangja{使われないデータが\\ラインを占める}{unused data\\occupies lines}};
  \draw[->, ln-muted, >={Stealth[length=1.4mm]}] (1.4,1.87) -- (2.9,1.87);
  \draw[->, ln-muted, >={Stealth[length=1.4mm]}] (4.6,1.87) -- (6.2,1.87);
\end{tikzpicture}

  \caption{4~KB と 256~KB のキャッシュについての，ブロックサイズの関数と
    してのミス率（形状のみ）．点は最良のブロックサイズを示す．}
  \label{fig:l14-block-size}
\end{figure}

\section{プリフェッチ}
\label{sec:l14-prefetching}

ブロックをプログラムが要求する前にキャッシュへ持ち込めば，ミスをそもそも
回避できる．\source{Lesson 14，動画 23--26；H\&P §2.3．}

\begin{definition}[プリフェッチ]
  近いうちにアクセスされると予想されるブロックを，アクセスされる前に
  キャッシュへ読み込むことを\term[ぷりふぇっち]{プリフェッチ}[prefetching]%
  という．
\end{definition}

プリフェッチは推測であり，その効果は推測が当たるかどうかで決まる．
\begin{itemize}
  \item 推測が当たれば，そのブロックへの最初のアクセスはヒットになる．
    メモリからの取得自体は行われるが，プロセッサが他の仕事をしている間に
    前倒しで行われるので，そのアクセスがミスペナルティを払うことはない．
  \item 推測が外れれば，その取得は無駄になる．さらに悪いことに，
    プリフェッチされたブロックは別のブロックの場所を奪っており，その
    ブロックが後で必要になればミスする．外れた推測はキャッシュを\emph{汚染}
    する．
\end{itemize}
プリフェッチが割に合うのは，推測の大半が当たる場合だけである．

\subsection{プリフェッチ命令}

プログラムは，次にどのデータにアクセスするかを知っていることが多い．
\term[ぷりふぇっちめいれい]{プリフェッチ命令}[prefetch instruction]とは，
与えられたアドレスを含むブロックをキャッシュへ持ち込むようプロセッサに
求める ISA の命令であり，データは変更しない．コンパイラやプログラマは，
あらかじめ分かっているアクセスの手前に，典型的には配列を巡るループの中に
これを置く \cite{callahan1991}．
\begin{lstlisting}[language=C, numbers=none]
for (i = 0; i < n; i++) {
  prefetch(&a[i + D]);    /* D 反復先 */
  sum = sum + a[i];
}
\end{lstlisting}
難しいのはプリフェッチ距離~$D$ である．
\begin{itemize}
  \item $D$ が小さすぎると，ループが要素 $i+D$ に達したときにプリフェッチが
    完了しておらず，そのアクセスは結局ブロックを待つので，得るものは少ない．
  \item $D$ が大きすぎると，ブロックは使われるずっと前に到着する．その間
    ラインを占有し，必要になる前に再び追い出されるかもしれない．
\end{itemize}
適切な距離は，取得にかかる時間を1回の反復の時間で割ったもの程度である．
どちらもマシンに依存するので，あるプロセッサに合う距離でコンパイルされた
プログラムは，別のプロセッサではプリフェッチが遅すぎたり早すぎたりする．
\caution{プリフェッチ距離は1つのマシンに合わせて調整されている．同じ
メモリでプロセッサが速くなれば，より大きな距離が必要になる．}

\begin{example}
  \label{ex:l14-prefetch-distance}
  プリフェッチされたブロックは，プリフェッチ命令の \SI{50}{\nano\second}
  後にキャッシュに入り，上のループの1回の反復には \SI{4}{\nano\second}
  かかる．距離は少なくとも $50/4 = 12.5$ でなければならないので，距離 16
  ならわずかな余裕がある．同じメモリで1回の反復を \SI{2}{\nano\second} で
  実行するプロセッサでは，ループが要素 $i+16$ に達するのはそのプリフェッチの
  \SI{32}{\nano\second} 後であり，ブロックの到着より \SI{18}{\nano\second}
  早い．このときループはメモリを待たなければならない．およそ 16 回に1回の
  アクセスが \SI{18}{\nano\second} 待つので，16 回の反復には
  \SI{32}{\nano\second} ではなく約 \SI{50}{\nano\second} かかり，速いはずの
  プロセッサはこのループを1反復当たり約 \SI{3.1}{\nano\second} でしか
  実行できない．少なくとも 25 の距離が必要である．
\end{example}

\subsection{ハードウェアプリフェッチ}

\term[はーどうぇあぷりふぇっち]{ハードウェアプリフェッチ}[hardware prefetching]%
では，プロセッサ自身がメモリアクセスの列を観測し，どのブロックが必要に
なるかを推測する．プログラムを変更する必要はなく，推測はプログラムが走る
マシンに適応する．3つの設計が異なるアクセスパターンを利用する
（\cref{fig:l14-prefetchers}）．
\begin{description}
  \item[ストリームバッファ] \term[すとりーむばっふぁ]{ストリームバッファ}[stream buffer]
    \cite{jouppi1990} は逐次的なプリフェッチャである．ブロックが順に
    アクセスされると，続く数ブロックをキャッシュの脇の小さなバッファに
    取得する．アクセスはメモリアクセスなしにそこからデータを得る．そして
    バッファが消費されるにつれて，さらに先のブロックを取得する．
  \item[ストライドプリフェッチャ] \term[すとらいどぷりふぇっちゃ]{ストライドプリフェッチャ}[stride prefetcher]%
    は，ループが配列の要素を $k$ 個ごとに訪れる場合のように，一定の距離
    ---\emph{ストライド}---で進むアクセスを検出し，1つ以上のストライド先の
    アドレスをプリフェッチする．
  \item[相関プリフェッチャ] \term[そうかんぷりふぇっちゃ]{相関プリフェッチャ}[correlating prefetcher]%
    は，あるブロックについて，その次にどのブロックがアクセスされたかを
    記録し，そのブロックが再びアクセスされたときに，記録された後続ブロックを
    プリフェッチする．メモリ中に散在するノードをもち，繰り返し辿られる
    連結リストのような，規則的な間隔をもたないパターンを扱える．
\end{description}

\begin{figure}[htbp]
  \centering
  % Access patterns exploited by three hardware prefetchers.  Numbers give
% the order of the accesses; P marks the block that is prefetched.
\begin{tikzpicture}[x=1cm,y=1cm, font=\footnotesize\sffamily, >={Stealth[length=1.5mm]}]
  \def\lnw{0.58}
  % \lnrow{y}{label}
  \def\lnrow#1#2{%
    \node[anchor=east, align=right] at (3.2,#1-0.29) {#2};
    \foreach \b in {0,...,11} {
      \draw[ln-rule, fill=white] ({3.4+\b*\lnw},#1) rectangle ++(\lnw,-0.58);
    }
  }
  % \lnacc{y}{block}{order}  accessed block
  \def\lnacc#1#2#3{%
    \draw[fill=ln-accent!25] ({3.4+#2*\lnw},#1) rectangle ++(\lnw,-0.58);
    \node at ({3.4+(#2+0.5)*\lnw},#1-0.29) {#3};
  }
  % \lnpre{y}{block}  prefetched block
  \def\lnpre#1#2{%
    \draw[fill=ln-tip!30] ({3.4+#2*\lnw},#1) rectangle ++(\lnw,-0.58);
    \node at ({3.4+(#2+0.5)*\lnw},#1-0.29) {P};
  }
  % block numbers
  \foreach \b in {0,...,11} {
    \node[font=\scriptsize\sffamily, text=ln-muted] at ({3.4+(\b+0.5)*\lnw},0.3) {\b};
  }
  \node[anchor=east, font=\scriptsize\sffamily, text=ln-muted] at (3.2,0.3) {\iflangja{ブロック}{block}};
  % (1) stream buffer: sequential
  \lnrow{0}{\iflangja{ストリームバッファ}{stream buffer}}
  \lnacc{0}{2}{1}\lnacc{0}{3}{2}\lnacc{0}{4}{3}
  \lnpre{0}{5}\lnpre{0}{6}
  \node[anchor=west, font=\scriptsize\sffamily, text=ln-muted] at (10.45,-0.29) {\iflangja{次のブロック}{next blocks}};
  % (2) stride prefetcher
  \lnrow{-1.0}{\iflangja{ストライド\\プリフェッチャ}{stride\\prefetcher}}
  \lnacc{-1.0}{1}{1}\lnacc{-1.0}{4}{2}\lnacc{-1.0}{7}{3}
  \lnpre{-1.0}{10}
  \node[anchor=west, font=\scriptsize\sffamily, text=ln-muted] at (10.45,-1.29) {\iflangja{ストライド 3}{stride 3}};
  % (3) correlating prefetcher
  \lnrow{-2.0}{\iflangja{相関\\プリフェッチャ}{correlating\\prefetcher}}
  \lnacc{-2.0}{8}{1}\lnacc{-2.0}{2}{2}\lnacc{-2.0}{5}{3}
  \draw[->, ln-accent] ({3.4+8.5*\lnw},-2.62) to[bend left=28] ({3.4+2.5*\lnw},-2.62);
  \draw[->, ln-accent] ({3.4+2.5*\lnw},-2.62) to[bend right=22] ({3.4+5.5*\lnw},-2.62);
  \node[anchor=west, font=\scriptsize\sffamily, text=ln-muted, align=left] at (10.45,-2.29)
    {\iflangja{8$\to$2，2$\to$5 を\\記録}{records\\8$\to$2, 2$\to$5}};
  % legend
  \draw[fill=ln-accent!25] (3.4,-3.55) rectangle ++(0.3,-0.3);
  \node[anchor=west, font=\scriptsize\sffamily] at (3.75,-3.7) {\iflangja{アクセス（数字は順序）}{accessed (number = order)}};
  \draw[fill=ln-tip!30] (7.4,-3.55) rectangle ++(0.3,-0.3);
  \node[anchor=west, font=\scriptsize\sffamily] at (7.75,-3.7) {\iflangja{プリフェッチ}{prefetched}};
\end{tikzpicture}

  \caption{3つのハードウェアプリフェッチャが利用するアクセスパターン．
    図示した3回のアクセスの後，ストリームバッファは続くブロックを，
    ストライドプリフェッチャは1ストライド先のブロックを取得し，相関
    プリフェッチャは次にブロック~8 がアクセスされたときにブロック~2 を
    取得する．}
  \label{fig:l14-prefetchers}
\end{figure}

\section{ループ交換}
\label{sec:l14-loop-interchange}

コンパイラは，プログラムがデータにアクセスする順序を変えることによっても
ミス率を下げられる．\source{Lesson 14，動画 27；H\&P §2.3．}C は2次元配列を
行ごとに格納する．\texttt{m[j][i]} と \texttt{m[j][i+1]} は隣接するが，
\texttt{m[j+1][i]} は1行分先にある．次のループの入れ子は，配列を列ごとに
走査する．
\begin{lstlisting}[language=C, numbers=none]
double m[2000][1000];
for (i = 0; i < 1000; i++)
  for (j = 0; j < 2000; j++)
    m[j][i] = 2 * m[j][i];
\end{lstlisting}
内側のループのアクセスは，どれも直前のアクセスから \num{1000} 要素，
\num{8000} バイト先へ飛ぶ．空間的局所性はまったくない．内側のループの
アクセスはどれも異なるブロックに触れ，内側のループの1巡が訪れる 2000 個の
ブロックは L1 キャッシュが保持できる数をはるかに超えるので，ほとんど
すべてのアクセスがミスする．

\begin{definition}[ループ交換]
  入れ子になったループの順序を入れ替え，最も内側のループが配列の最後の
  添字を走ってデータが逐次的にアクセスされるようにするコンパイラ変換を
  \term[るーぷこうかん]{ループ交換}[loop interchange]という．
\end{definition}

交換の後，外側のループは \texttt{j} を，内側のループは \texttt{i} を走る．
内側のループは隣接する要素を訪れるようになり，1つのブロックが連続する
複数のアクセスに応え，逐次プリフェッチャやストライドプリフェッチャが次の
ブロックを予測できる．コンパイラがループを交換してよいのは，結果が
変わらないことを証明できる場合だけである．ここでは各反復が自分の要素だけを
読み書きするので，反復はどの順序で実行してもよい．

\begin{keypoint}
  プリフェッチはブロックへの最初のアクセスのミスさえ回避できるが，役に
  立つのは推測が正しく，かつ間に合う場合だけである．ループ交換は，アクセスの
  順序をメモリ上のデータの配置に合わせることで，プログラムの側からミスを
  攻める．
\end{keypoint}

\section{ミスの重なり}
\label{sec:l14-overlap}

ミスペナルティは，ブロックの取得に要する時間より大きくなりうる．
\source{Lesson 14，動画 28--30．}アウトオブオーダプロセッサ（第7章）は
多くの命令を同時に処理しており，そのロードとストアのいくつかは数サイクルの
うちにメモリへアクセスしうる．欠けているブロックの取得には多くのサイクルが
かかり，その間に別のアクセスもミスしうる．

\begin{definition}[オーバーラップミス]
  先行するミスがまだ処理されている間に生じるキャッシュミスを
  \term[おーばーらっぷみす]{オーバーラップミス}[overlap miss]という．
\end{definition}

ロードがミスした後も，アウトオブオーダプロセッサは他の命令を実行し続ける．
しかしそのロードはコミットできないので，プロセッサはやがて ROB エントリか
リザベーションステーションを使い果たして停止する．その時点より前に発行
された2つ目のミスは，1つ目と重なる．それがどれだけの代償になるかは
キャッシュによる．
\begin{itemize}
  \item \term[ぶろっきんぐきゃっしゅ]{ブロッキングキャッシュ}[blocking cache]%
    は，ミスを処理している間は他に何もしない．後続のアクセスは先行する
    ミスが完了するまで探索さえされないので，重なったミスは自身のペナルティに
    加えて先行するミスの残り時間を払う（\cref{fig:l14-overlap}a）．
  \item \term[のんぶろっきんぐきゃっしゅ]{ノンブロッキングキャッシュ}[non-blocking cache]%
    は，ミスの処理中にもヒットに応え（\emph{ヒットアンダーミス}），さらに
    別のミスを開始することもできる．このとき複数の要求がメモリで同時に
    処理され，メモリが要求を並列に処理できる限り，\emph{メモリレベル並列性}が
    活用される．
\end{itemize}

\begin{definition}[ミスアンダーミス対応]
  \term[みすあんだーみすたいおう]{ミスアンダーミス対応}[miss-under-miss support]%
  をもつキャッシュは，複数のミスを同時に処理できる．処理中の各ミスを記録し，
  それらの要求をメモリ階層の次のレベルへ並行して送る．
\end{definition}

処理中の各ミスは\term[みすじょうたいかんりれじすた]{ミス状態管理レジスタ}[miss status handling register]%
\index{MSHR|see{ミス状態管理レジスタ}}（MSHR）に記録される．キャッシュで
ミスしたアクセスは，まず自身のブロックアドレスを MSHR と比較する．
\begin{description}
  \item[一致なし] そのアクセスは新しいミスである．空いている MSHR を1つ
    確保し，そこにブロックと起こすべき命令を記録して，要求をメモリへ送る．
  \item[一致あり] そのブロックはすでに到着の途中である．これを
    \term[はーふみす]{ハーフミス}[half-miss]という．2つ目の要求は送られず，
    命令は一致した MSHR に追加されるだけである．
\end{description}
データが到着すると，それがキャッシュに置かれ，その MSHR に記録された
すべての命令が起こされ，MSHR が解放される．MSHR が2つあるだけでも
ブロッキングキャッシュより大きく改善し，4つあればさらに良い．そして
メモリアクセスは非常に長くかかるので，16〜32 個の MSHR でも効果が続く．
\margin{実際の機械が追跡できる数はもっと少ない．Nehalem から Skylake まで
の Intel のコアはラインフィルバッファを 10 個もち，コア当たり同時に処理
できる L1 データミスは 10 件までである．Intel 最適化マニュアル．}

\begin{figure}[htbp]
  \centering
  % Two overlapping misses with a miss penalty of 20 cycles; the second miss
% occurs 6 cycles after the first.  (a) The cache services one miss at a
% time (blocking cache); (b) the cache has miss-under-miss support.  The red line marks the
% moment the second miss occurs.
\begin{tikzpicture}[x=0.215cm,y=1cm, font=\footnotesize\sffamily,
  svc/.style={draw, fill=ln-accent!25, minimum height=0.42cm, inner sep=0pt, font=\scriptsize\sffamily},
  lab/.style={anchor=east, font=\footnotesize\sffamily}]
  % \lnbar{x0}{x1}{y}{style}{text}
  \def\lnbar#1#2#3#4#5{%
    \draw[#4] (#1,#3+0.21) rectangle (#2,#3-0.21);
    \node[font=\scriptsize\sffamily] at ({(#1+#2)/2},#3) {#5};
  }
  % ---- (a) one miss at a time
  \node[anchor=west] at (-5,0.75) {\iflangja{(a) ブロッキングキャッシュ}{(a) blocking cache}};
  \node[lab] at (-1,0) {\iflangja{ミス 1}{miss 1}};
  \lnbar{0}{20}{0}{fill=ln-accent!25}{\iflangja{処理}{service}}
  \node[lab] at (-1,-0.6) {\iflangja{ミス 2}{miss 2}};
  \lnbar{6}{20}{-0.6}{fill=ln-box, draw=ln-rule}{\iflangja{待ち}{wait}}
  \lnbar{20}{40}{-0.6}{fill=ln-accent!25}{\iflangja{処理}{service}}
  \draw[ln-caution, thick] (6,-0.25) -- (6,-0.95);
  % ---- (b) miss-under-miss support
  \node[anchor=west] at (-5,-1.45) {\iflangja{(b) ミスアンダーミス対応}{(b) miss-under-miss support}};
  \node[lab] at (-1,-2.2) {\iflangja{ミス 1}{miss 1}};
  \lnbar{0}{20}{-2.2}{fill=ln-accent!25}{\iflangja{処理}{service}}
  \node[lab] at (-1,-2.8) {\iflangja{ミス 2}{miss 2}};
  \lnbar{6}{26}{-2.8}{fill=ln-accent!25}{\iflangja{処理}{service}}
  \draw[ln-caution, thick] (6,-2.45) -- (6,-3.15);
  % ---- time axis
  \draw[->] (0,-3.45) -- (42.5,-3.45);
  \foreach \t in {0,10,20,30,40} {
    \draw (\t,-3.45) -- (\t,-3.53) node[below, font=\scriptsize\sffamily] {\t};
  }
  \node[anchor=west, font=\scriptsize\sffamily] at (43,-3.45) {\iflangja{サイクル}{cycles}};
\end{tikzpicture}

  \caption{ミスペナルティが 20 サイクルの2つのミス．2つ目は1つ目の6
    サイクル後に生じる（縦線）．(a)~ブロッキングキャッシュでは，2つ目の
    ミスに 20 サイクルではなく 34 サイクルかかる．(b)~ミスアンダーミス
    対応があれば，どちらのミスも 20 サイクルで済む．}
  \label{fig:l14-overlap}
\end{figure}

ミスアンダーミス対応には追加のハードウェアが必要である．多くのメモリ
アクセスを同時に発行するプロセッサ，とりわけアウトオブオーダプロセッサで
割に合う．

\section{キャッシュ階層}
\label{sec:l14-hierarchy}

L1 キャッシュのミスペナルティは，主記憶のアクセス時間に支配される．
\source{Lesson 14，動画 31--34；H\&P §B.3．}これは，そもそも L1 キャッシュを
生んだのと同じ発想で下げられる．すなわち，L1 と主記憶の間にもう1つ
キャッシュを置く．L1 でミスしたアクセスは
\term[L2 きゃっしゅ]{L2 キャッシュ}[L2 cache]を探索する．L2 は L1 より
大きく遅いが，それでも主記憶よりはるかに速い．そして L2 でミスしたものだけが
さらに先へ進む．

\begin{definition}[キャッシュ階層]
  \term[きゃっしゅかいそう]{キャッシュ階層}[cache hierarchy]とは，L1，L2，
  L3，\ldots という複数のレベルのキャッシュからなり，各レベルは直前の
  レベルがミスしたときにだけ探索されるものである
  （\cref{fig:l14-hierarchy}）．主記憶の直前のレベルを
  \term[さいしゅうれべるきゃっしゅ]{最終レベルキャッシュ}[last-level cache]%
  \index{LLC|see{最終レベルキャッシュ}}（LLC）という．
\end{definition}

\begin{figure}[htbp]
  \centering
  % A three-level cache hierarchy: each level is consulted only on a miss in
% the level before it; the miss penalty of a level is the average access
% time of everything behind it.
\begin{tikzpicture}[x=1cm,y=1cm, font=\footnotesize\sffamily, >={Stealth[length=1.8mm]},
  lvl/.style={draw, fill=ln-accent!15, align=center, inner sep=2pt},
  lab/.style={font=\scriptsize\sffamily, text=ln-muted}]
  \node[draw, fill=ln-box, minimum width=1.5cm, minimum height=0.9cm, align=center] (cpu) at (0,0)
    {\iflangja{プロセッサ}{processor}};
  \node[lvl, minimum width=1.0cm, minimum height=0.7cm] (l1) at (2.2,0) {L1};
  \node[lvl, minimum width=1.4cm, minimum height=1.0cm] (l2) at (4.3,0) {L2};
  \node[lvl, minimum width=1.9cm, minimum height=1.3cm] (l3) at (6.9,0) {L3\\(LLC)};
  \node[draw, fill=ln-box, minimum width=2.1cm, minimum height=1.6cm, align=center] (mm) at (10.0,0)
    {\iflangja{主記憶}{main\\memory}};
  \draw[->] (cpu) -- (l1);
  \draw[->] (l1) -- node[lab, above] {\iflangja{ミス}{miss}} (l2);
  \draw[->] (l2) -- node[lab, above] {\iflangja{ミス}{miss}} (l3);
  \draw[->] (l3) -- node[lab, above] {\iflangja{ミス}{miss}} (mm);
  % miss penalties
  \draw[ln-muted] (5.95,-0.9) -- (5.95,-1.05) -- (11.05,-1.05) -- (11.05,-0.9);
  \node[lab, anchor=north] at (8.5,-1.05) {\iflangja{L2 のミスペナルティ}{miss penalty of L2}};
  \draw[ln-muted] (3.6,-1.6) -- (3.6,-1.75) -- (11.05,-1.75) -- (11.05,-1.6);
  \node[lab, anchor=north] at (7.3,-1.75) {\iflangja{L1 のミスペナルティ}{miss penalty of L1}};
  % size / speed trend
  \draw[<->, ln-muted] (1.6,1.15) -- (11.0,1.15);
  \node[lab, anchor=south west] at (1.6,1.18) {\iflangja{小さく高速}{smaller, faster}};
  \node[lab, anchor=south east] at (11.0,1.18) {\iflangja{大きく低速}{larger, slower}};
\end{tikzpicture}

  \caption{3レベルのキャッシュ階層．各レベルのミスペナルティは，その後ろに
    あるレベル群の平均アクセス時間である．}
  \label{fig:l14-hierarchy}
\end{figure}

\subsection{キャッシュ階層の AMAT}

あるレベルのミスペナルティは，その後ろのレベルからブロックを得るのに
かかる平均時間，すなわち階層の残りの部分の AMAT である．
\cref{eq:l14-amat} をレベルごとに適用すると次が得られる．
\begin{equation}
  \begin{aligned}
    \text{AMAT} &= h_1 + m_1 \times P_1 , \\
    P_1 &= h_2 + m_2 \times P_2 , \\
    &\;\;\vdots \\
    P_{\text{LLC}} &= \text{主記憶のアクセス時間},
  \end{aligned}
  \label{eq:l14-amat-hierarchy}
\end{equation}
ここで $h_i$ はレベル~$i$ のヒット時間，$m_i$ はミス率，$P_i$ はミス
ペナルティである．ミス率 $m_i$ は，レベル~$i$ に到達したアクセスについて
数える（\cref{sec:l14-local-global}）．

\subsection{L1 と L2}

\cref{eq:l14-amat-hierarchy} では，L1 のヒット時間 $h_1$ はすべての
アクセスが払うが，L2 のヒット時間 $h_2$ は L1 でミスしたアクセスの割合
$m_1$ だけが払う．したがって L1 は，前の節までの技法によって短いヒット
時間を目指して設計され，L2 はヒット時間が長くなる代わりにミス率を下げて
よい．\margin{Skylake の値（Intel 最適化マニュアル）は，L1 データ 32~KB で
4 サイクル，L2 256~KB で 12 サイクル，L3 はコア当たり約 2~MB で 44
サイクル，主記憶は 60--100~ns，すなわち数百サイクルである．}

\begin{example}
  \label{ex:l14-hierarchy}
  主記憶は1回のアクセスに \SI{50}{\nano\second} かかる．32~KB の
  キャッシュはヒット時間 \SI{1}{\nano\second}，ミス率 \SI{10}{\percent}，
  256~KB のキャッシュはヒット時間 \SI{5}{\nano\second}，ミス率
  \SI{3}{\percent} である．階層の中でも 256~KB のキャッシュは単独の場合と
  同じアクセスでミスする，すなわちそこに到達するアクセスの
  $0.03/0.10 = \SI{30}{\percent}$ でミスすると仮定する．
  \begin{center}\small
    \begin{tabular}{@{}ll@{}}
      \toprule
      構成 & AMAT \\
      \midrule
      キャッシュなし                & \SI{50}{\nano\second} \\
      32~KB キャッシュのみ          & $1 + 0.10\times 50 = \SI{6}{\nano\second}$ \\
      256~KB キャッシュのみ         & $5 + 0.03\times 50 = \SI{6.5}{\nano\second}$ \\
      L1 32~KB，L2 256~KB           & $1 + 0.10\times(5 + 0.30\times 50) = \SI{3}{\nano\second}$ \\
      \bottomrule
    \end{tabular}
  \end{center}
  この階層は，どちらのキャッシュ単独よりも2倍速い．すべてのアクセスが払う
  小さなキャッシュの短いヒット時間と，長いヒット時間をアクセスの
  \SI{10}{\percent} だけが払う大きなキャッシュの低いミス率とを兼ね備えて
  いる．
\end{example}

\section{局所ミス率と大域ミス率}
\label{sec:l14-local-global}

\cref{ex:l14-hierarchy} では，L2 キャッシュはそこに到達するアクセスの
\SI{30}{\percent} でミスし，これは L1 の3倍の率である．\source{Lesson 14，
動画 35--37；H\&P §B.3．}だからといって L2 が劣ったキャッシュだという
ことにはならない．L2 が見るのは L1 がミスしたアクセスだけであり，それらは
局所性が最も弱いものである．応えるのが容易なアクセスは L1 より先へは
行かない．したがって下位レベルのミス率は，その基準とともに述べなければ
ならない．

\begin{definition}[局所率と大域率]
  キャッシュの\term[きょくしょひっとりつ]{局所ヒット率}[local hit rate]と
  は，そのヒット数を，そこに到達したアクセス数で割ったものであり，
  \term[きょくしょみすりつ]{局所ミス率}[local miss rate]は
  $1 - \text{局所ヒット率}$ である．

  キャッシュの\term[たいいきみすりつ]{大域ミス率}[global miss rate]とは，
  そのミス数を，プロセッサのすべてのメモリアクセス数で割ったものであり，
  \term[たいいきひっとりつ]{大域ヒット率}[global hit rate]は
  $1 - \text{大域ミス率}$ である．
\end{definition}

L1 では局所率と大域率は一致する．下位のレベルでは，そこに到達するアクセスは
上のレベルのミスなので，レベル~$n$ の大域ミス率は局所ミス率の積になる．
\begin{equation}
  \text{大域ミス率}_n = \prod_{i=1}^{n} \text{局所ミス率}_i .
  \label{eq:l14-global-miss}
\end{equation}
\cref{eq:l14-amat-hierarchy} のミス率は局所ミス率である．あるレベルの
大域ミス率は，全アクセスのうちどれだけの割合がそのレベルより先へ進まなければ
ならないかを述べるので，レベルや階層を比較するのに適した尺度である．

\begin{definition}[1000 命令当たりミス数]
  \term[1000めいれいあたりみすすう]{1000 命令当たりミス数}[misses per 1000 instructions]%
  \index{MPKI|see{1000 命令当たりミス数}}（MPKI）という指標は，実行された
  1000 命令当たりのキャッシュのミス数を数える．
\end{definition}

大域ミス率と同様に MPKI もあるレベルのすべてのミスを数えるが，基準が
メモリアクセスではなく命令である．
\begin{equation}
  \text{MPKI} = 1000 \times \text{大域ミス率} \times
  \frac{\text{メモリアクセス数}}{\text{命令}} .
  \label{eq:l14-mpki}
\end{equation}
これにより，そのレベルのミスペナルティをプログラムが命令当たりどれだけの
頻度で払うかが直接に分かる．\tip{典型的なコードでは命令の約3分の1がメモリ
にアクセスするので，MPKI $\approx 330 \times$ 大域ミス率 である（H\&P の
命令構成比）．}

\begin{example}
  \label{ex:l14-local-global}
  あるプログラムが \num{1000000} 命令を実行し，メモリアクセスは
  \num{1500000} 回である．L1 は \num{90000} 回，L2 は \num{27000} 回
  ミスする．
  \begin{itemize}
    \item L1: 局所ミス率と大域ミス率は
      $\num{90000}/\num{1500000} = \SI{6}{\percent}$，MPKI は
      $\num{90000}/1000 = 90$．
    \item L2: 局所ミス率は $\num{27000}/\num{90000} = \SI{30}{\percent}$
      （局所ヒット率 \SI{70}{\percent}），大域ミス率は
      $\num{27000}/\num{1500000} = \SI{1.8}{\percent} = 0.06\times 0.30$
      （大域ヒット率 \SI{98.2}{\percent}），MPKI は $27$．
  \end{itemize}
\end{example}

\begin{keypoint}
  AMAT は局所ミス率で計算する．レベルや階層の比較には大域ミス率か MPKI を
  用いる．上のレベルのミスだけを見るレベルでは，局所ヒット率が低いのは
  当然だからである．
\end{keypoint}

\section{包含性}
\label{sec:l14-inclusion}

L1 にあるブロックは，L2 にもあるかもしれないし，ないかもしれない．
\source{Lesson 14，動画 38--40；H\&P §B.3；\cite{baer1988}．}レベル間の
関係には3つの可能性がある．
\begin{description}
  \item[包含] L1 にあるブロックはすべて L2 にもある．
  \item[排他] L1 にあるブロックは L2 にはない．
  \item[どちらでもない] L1 にあるブロックの一部は L2 にもあり，他はない．
\end{description}

\begin{definition}[包含性]
  あるレベルにあるブロックがすべて，次の，より大きなレベルにも存在するとき，
  そのキャッシュ階層は\term[ほうがんせい]{包含性}[inclusion property]を
  もつという．
\end{definition}

LRU 置換をもつ大きな L2 なら，L1 のブロックを自動的にすべて保持しそうに
思える．そうではない．L1 でのヒットは L2 に届かないので，頻繁に使われ，
そのため L1 に留まるブロックは，L2 の LRU カウンタからは使われていないように
見える．やがてそれは L2 で最も長く使われていないブロックとなり，L1 に
残ったまま L2 では追い出される．

\begin{example}
  \label{ex:l14-inclusion}
  L1 は 2 ブロック，L2 は 3 ブロックを保持する．どちらもフルアソシアティブで
  LRU 置換をもち，初期状態は空であり，両方のレベルでミスしたブロックは
  両方に読み込まれる．\cref{tab:l14-inclusion} はアクセス P，Q，P，R，P，S
  を追う．P への後のアクセスは L1 でヒットするので，L2 が P を最後に見るのは
  最初のアクセスである．S でのミスの時点で P は L2 で最も長く使われていない
  ブロックであり，まだ L1 にあるにもかかわらず追い出される．
\end{example}

\begin{table}[htbp]
  \centering\small
  \caption{両方のレベルが LRU 置換をもつ2レベルの階層
    （\cref{ex:l14-inclusion}）．内容は，最も長く使われていないものから
    最も最近使われたものの順に並べる．最後のアクセスで包含性が破れる．}
  \label{tab:l14-inclusion}
  \begin{tabular}{@{}lllll@{}}
    \toprule
    アクセス & L1   & L2   & 後の L1 & 後の L2 \\
    \midrule
    P      & ミス   & ミス & P        & P \\
    Q      & ミス   & ミス & P Q      & P Q \\
    P      & ヒット & ---  & Q P      & P Q \\
    R      & ミス   & ミス & P R      & P Q R \\
    P      & ヒット & ---  & R P      & P Q R \\
    S      & ミス   & ミス & P S      & Q R S \\
    \bottomrule
  \end{tabular}
\end{table}

したがって包含性が欲しければ，それを強制しなければならない．1つの方法は，
L2 のすべてのラインに\emph{包含ビット}を設け，そのブロックが L1 にもある間
これを立てておくことである．L2 は，このビットが立っているラインを犠牲に
選ばない．\sidenote{どちらも実際に使われている．Skylake までの Intel の
クライアント向け L3 は包含的であり，フィルタも兼ねる．L3 にないブロックは
どのコアのキャッシュにもないので，他のコアに問い合わせる必要がない
（第19章）．Skylake-SP と AMD Zen の L3 は非包含であり，主として上位から
追い出されたブロックを受け取るので，容量が写しに費やされない．}

この章で，単一プロセッサにおけるキャッシュの議論は完結する（第19章では
マルチプロセッサのキャッシュに戻る）．第15章では，キャッシュと主記憶を
作るメモリ技術に話を移す．

\needspace{10\baselineskip}
\begin{keypoint}[章のまとめ]
  AMAT $=$ ヒット時間 $+$ ミス率 $\times$ ミスペナルティ であり，各項には
  それぞれの技法がある．ヒット時間は，探索のパイプライン化，仮想アドレスで
  インデックスしながら TLB が変換し，
  $\text{キャッシュサイズ} \le N \times \text{ページサイズ}$ ならエイリアシングを
  避けられる VIPT キャッシュ，ウェイ予測，そして安価な置換アルゴリズムで
  ある NMRU と PLRU によって短くなる．ミスは初期参照ミス，容量性ミス，
  競合性ミスのいずれかであり，ミス率は適切なブロックサイズ，プリフェッチ
  （プリフェッチ命令，ストリームバッファ，ストライドプリフェッチャ，相関
  プリフェッチャ），ループ交換によって下がる．ミスペナルティは，処理中の
  ミスを MSHR で追跡するミスアンダーミス対応と，レベルごとに局所ミス率で
  AMAT を計算するキャッシュ階層によって下がる．大域ミス率と MPKI は，
  あるレベルより先へ進むアクセスがどれだけあるかを表す．各レベルが LRU を
  用いる場合，包含性は自動的には成り立たない．
\end{keypoint}

\end{document}
